\chapter{Cấu hình và Secret trong cụm}
\label{ch:secrets}

\section{Hai cơ chế tiêm cấu hình}

Phần~I đã xác lập rằng biến môi trường thắng mọi file. Kubernetes cung
cấp hai cách để đặt chúng, và các manifest dùng cả hai một cách có chủ
đích:

\begin{lstlisting}[language=yaml]
env:                                     # literal, visible in git
  - { name: SPRING_DATASOURCE_URL,
      value: "jdbc:postgresql://postgres-course:5432/ts_course" }
envFrom:                                 # every key of a Secret
  - secretRef: { name: postgres-course-app }
  - secretRef: { name: jwt }
\end{lstlisting}

Chính sự tách bạch này \emph{là} ranh giới bảo mật: topology (hostname,
cổng, feature flag) không phải bí mật và thuộc về git, nơi có thể review;
thông tin đăng nhập (credential) không bao giờ xuất hiện trong git. Khi
đọc một Deployment ở đây, mọi mục \code{env:} là cấu hình công khai và
mọi dòng \code{envFrom:} là con trỏ tới thứ được tạo ngoài luồng.

\section{Secret: là gì và không là gì}

Secret là một object thuộc namespace, chứa các cặp key/value được
encode base64. Hai điều thẳng thắn về nó:

\begin{itemize}
  \item Base64 là \emph{encoding}, không phải mã hóa.
  \code{echo cm9vdA== | base64 -d} lộ ngay mật khẩu. Thứ bảo vệ nó là
  kiểm soát truy cập (RBAC quyết định ai được đọc Secret --- để ý
  Chương~22 không cấp cho Jenkins verb nào như vậy), vòng đời tách khỏi
  code, và tùy chọn mã hóa khi lưu (encryption at rest).
  \item Nó vẫn là chỗ chứa đúng: mọi cơ chế cao hơn nó (Sealed Secrets,
  Vault, secret manager của cloud) cuối cùng đều vật chất hóa thành một
  Secret để pod tiêu thụ. Học object thuần trước là học chính giao diện.
\end{itemize}

\section{\code{create-secrets.sh}, có chú giải}

Secret ở đây được tạo bằng một script lũy đẳng, không phải manifest:

\begin{lstlisting}[language=bash]
kubectl -n "$NS" create secret generic postgres-course \
  --from-literal=POSTGRES_USER="$PG_USER" \
  --from-literal=POSTGRES_PASSWORD="$PG_PASSWORD" \
  --from-literal=POSTGRES_DB=ts_course \
  --dry-run=client -o yaml | kubectl apply -f -     (*@\co{1}@*)

kubectl -n "$NS" create secret generic postgres-course-app \
  --from-literal=SPRING_DATASOURCE_USERNAME="$PG_USER" \
  --from-literal=SPRING_DATASOURCE_PASSWORD="$PG_PASSWORD" \
  --dry-run=client -o yaml | kubectl apply -f -     (*@\co{2}@*)

kubectl -n "$NS" create secret generic jwt \
  --from-literal=JWT_SECRET="$JWT_SECRET"           (*@\co{3}@*)
\end{lstlisting}

\begin{description}
  \item[\co{1}] Thành ngữ tạo-lũy-đẳng: \code{create} đơn thuần sẽ thất
  bại nếu Secret đã tồn tại; render bằng \code{--dry-run} rồi pipe sang
  \code{apply} khiến script chạy lại an toàn --- tạo mới hay cập nhật,
  đằng nào cũng hội tụ. Một bài học nhỏ về khai báo-hơn-mệnh-lệnh, dưới
  dạng shell.
  \item[\co{2}] Cùng một credential viết hai lần, theo cách mỗi bên tiêu
  thụ muốn: image Postgres cần \code{POSTGRES\_USER}; Spring cần
  \code{SPRING\_DATASOURCE\_USERNAME}. Hai Secret nghĩa là
  \code{envFrom} của mỗi pod chỉ nhập những tên nó hiểu --- không pod nào
  phải ôm một mớ biến lạ.
  \item[\co{3}] Một Secret được hai Deployment tiêu thụ --- auth-service
  dùng nó để ký JWT, api-gateway dùng để xác minh. Hợp đồng khóa-chia-sẻ
  của Chương~5, thể hiện dưới dạng hạ tầng: xoay vòng (rotate) khóa là
  một lần chạy script cộng hai lần khởi động lại pod, và hai bên tiêu thụ
  \emph{không thể} lệch nhau.
\end{description}

\begin{decision}[script, không phải manifest --- và không nằm trong kustomization]
Manifest Secret trong git sẽ đặt credential đã base64 chỉ cách bất kỳ ai
có quyền truy cập repo một lệnh \code{base64 -d}, mãi mãi, trong lịch
sử. Script giữ giá trị trong môi trường của người vận hành
(\code{JWT\_SECRET=... ./create-secrets.sh}), giá trị mặc định khớp với
dev nên hành vi giống hệt, và git chỉ giữ \emph{hình dạng}. Comment
trong kustomization nói cùng một điều từ phía bên kia: secret cố tình
vắng mặt trong \code{resources:}.
\end{decision}

\section{ConfigMap}

Người anh em không-bí-mật của Secret --- cùng cơ chế, không giả vờ nhạy
cảm, cộng thêm một tài năng mà Secret cũng có nhưng ConfigMap dùng nhiều
hơn: mount thành \emph{file}. Dự án này chưa có ConfigMap nào, vì một lý
do thẳng thắn: cấu hình không-bí-mật của nó đã chảy qua config-server và
biến môi trường. ConfigMap đầu tiên sẽ đến một cách tự nhiên cùng việc
cho config-server nghỉ hưu ở Chương~27 --- các file
\code{configurations/*.yml} được mount vào pod, để Spring đọc trực tiếp
và xóa hẳn một service.

\begin{pitfall}[đổi Secret rồi, chẳng có gì xảy ra]
Triệu chứng: bạn chạy lại \code{create-secrets.sh} với mật khẩu mới; các
service vẫn dùng mật khẩu cũ, cho tới khi vài ngày sau một pod nào đó
khởi động lại và bắt đầu lỗi xác thực --- một sự cố hẹn giờ.
Nguyên nhân: biến môi trường được tiêm \emph{lúc container khởi động};
pod đang chạy không bao giờ thấy Secret cập nhật. Vì vậy xoay vòng theo
thiết kế là hai bước: cập nhật Secret, rồi \code{kubectl -n ts rollout
restart deployment ...} cho mọi bên tiêu thụ, ngay lập tức --- trong lúc
credential cũ và mới đều còn dùng được, nếu hệ thống phía sau cho phép.
\end{pitfall}

\section{Ngoài phạm vi dự án}

Khoảng cách giữa thiết lập này và thực hành công nghiệp là một khái
niệm: làm cho secret \emph{giao được qua pipeline} mà không đọc được
trong git. Ba nấc trên chiếc thang đó: \textbf{Sealed Secrets} (mã hóa
bằng khóa công khai của cụm; commit khối đã niêm phong; một controller
giải mã trong cụm --- tương thích GitOps với một bộ phận chuyển động);
\textbf{External Secrets Operator} (một manifest Secret là
\emph{con trỏ} vào Vault/AWS Secrets Manager, được controller đồng bộ);
và \textbf{Vault injection} (secret không bao giờ trở thành object
Kubernetes). Mỗi cách đều thay thế script; không cách nào thay đổi cách
pod tiêu thụ --- \code{envFrom} sống sót qua mọi cuộc di cư, và đó là lý
do các manifest đã sẵn sàng cho tương lai.

\begin{quiz}
\begin{enumerate}
  \item Base64 không phải mã hóa --- vậy hãy kể ba thứ thực sự bảo vệ
  một Secret, và dự án này đứng ở đâu với từng thứ.
  \item Vì sao hai Secret cho một credential Postgres? Việc tách này
  ngăn được lỗi gì?
  \item Mô tả trọn vẹn việc xoay vòng JWT secret trên cụm này: lệnh,
  thứ tự, và kiểu lỗi khi chỉ làm một nửa.
  \item Mục cấu hình nào trong dự án sẽ chuyển sang ConfigMap trước
  tiên, và điều đó cho service nào nghỉ hưu?
\end{enumerate}
\end{quiz}
